\documentclass{beamer}

\usepackage[T2A]{fontenc}
\usepackage[russian]{babel}
\usepackage[utf8]{inputenc}
\usepackage{graphicx}
\usepackage{xspace}
\usepackage{psfrag}
\usepackage{hyperref}
\usepackage[normalem]{ulem}
\usepackage{booktabs}
\usepackage{tabularx}

\newcolumntype{Y}{>{\raggedright\arraybackslash}X}

\beamertemplatenavigationsymbolsempty
\setbeamertemplate{footline}[frame number]
\usecolortheme{seahorse}

\DeclareSymbolFont{greek}{U}{eur}{m}{n}                                                        
\SetSymbolFont{greek}{bold}{U}{eur}{b}{n}                                                      
\DeclareSymbolFontAlphabet{\gr}{greek}                                                         
\DeclareMathSymbol{\alpha}{\mathord}{greek}{"0B}                                               
\DeclareMathSymbol{\beta}{\mathord}{greek}{"0C}                                                
\DeclareMathSymbol{\gamma}{\mathord}{greek}{"0D}                                               
\DeclareMathSymbol{\delta}{\mathord}{greek}{"0E}                                               
\DeclareMathSymbol{\epsilon}{\mathord}{greek}{"0F}                                             
\DeclareMathSymbol{\zeta}{\mathord}{greek}{"10}                                                
\DeclareMathSymbol{\eta}{\mathord}{greek}{"11}                                                 
\DeclareMathSymbol{\theta}{\mathord}{greek}{"12}                                               
\DeclareMathSymbol{\iota}{\mathord}{greek}{"13}                                                
\DeclareMathSymbol{\kappa}{\mathord}{greek}{"14}                                               
\DeclareMathSymbol{\lambda}{\mathord}{greek}{"15}                                              
\DeclareMathSymbol{\mu}{\mathord}{greek}{"16}                                                  
\DeclareMathSymbol{\nu}{\mathord}{greek}{"17}                                                  
\DeclareMathSymbol{\xi}{\mathord}{greek}{"18}                                                  
\DeclareMathSymbol{\pi}{\mathord}{greek}{"19}                                                  
\DeclareMathSymbol{\rho}{\mathord}{greek}{"1A}                                                 
\DeclareMathSymbol{\sigma}{\mathord}{greek}{"1B}                                               
\DeclareMathSymbol{\tau}{\mathord}{greek}{"1C}                                                 
\DeclareMathSymbol{\upsilon}{\mathord}{greek}{"1D}                                             
\DeclareMathSymbol{\phi}{\mathord}{greek}{"1E}                                                 
\DeclareMathSymbol{\chi}{\mathord}{greek}{"1F}                                                 
\DeclareMathSymbol{\psi}{\mathord}{greek}{"20}                                                 
\DeclareMathSymbol{\omega}{\mathord}{greek}{"21}                                               
\DeclareMathSymbol{\varepsilon}{\mathord}{greek}{"22}                                          
\DeclareMathSymbol{\vartheta}{\mathord}{greek}{"23}                                            
\DeclareMathSymbol{\varpi}{\mathord}{greek}{"24}                                               
\let\varrho\rho                                                                                
\let\varsigma\sigma                                                                            
\DeclareMathSymbol{\varphi}{\mathord}{greek}{"27}

\title{Разработка датасета и сопутствующего инструментария
 для~математического вывода}
\author{Макар Гороховик}
\date{\today}

\begin{document}
\maketitle

% Слайд 1
\begin{frame}
  \frametitle{Введение: Проект SHARP}
  \begin{center}
      \Huge \textbf{SHARP} \\[0.5em]
      \Large Span-level Hallucination Annotation \\ for Reasoning Paths
  \end{center}
  \vspace{1em}
  
  \large Датасет для детекции галлюцинаций на уровне фрагментов в математических рассуждениях.
  
  \vspace{1em}
  \normalsize
  \begin{itemize}
      \item \textbf{11K} примеров
      \item \textbf{Span-level} аннотации
      \item \textbf{Open Source}
  \end{itemize}
\end{frame}

% Слайд 2
\begin{frame}
  \frametitle{Проблема: Галлюцинации в LLM}
  \footnotesize
  
  \begin{block}{Сильные стороны LLM}
      Современные большие языковые модели демонстрируют выдающиеся способности в решении сложных математических задач при генерации пошаговых решений chain-of-thought.
  \end{block}

  \begin{alertblock}{Критическая проблема}
      Модели остаются склонными к галлюцинациям — генерации правдоподобного, но логически некорректного или фактически неверного контента.
  \end{alertblock}

  \begin{block}{Эффект распространения}
      \begin{enumerate}
          \item Один ошибочный логический шаг $\rightarrow$
          \item Распространяется через всё решение $\rightarrow$
          \item Испорченный финальный ответ
      \end{enumerate}
  \end{block}
  \small \textit{Детекция таких ошибок критически важна для повышения надёжности LLM-генерации рассуждений.}
\end{frame}

% Слайд 3
\begin{frame}
  \frametitle{Существующие подходы к детекции ошибок}
  
  \begin{columns}[T]
      \begin{column}{0.48\textwidth}
          \begin{block}{ORM (Outcome Reward Models)}
              Присваивают награду на основе \textbf{финального ответа} траектории рассуждений. \\
              \vspace{1em}
              \footnotesize \textit{Ограничение: не позволяет идентифицировать конкретный ошибочный шаг в цепочке рассуждений.}
          \end{block}
      \end{column}
      
      \begin{column}{0.48\textwidth}
          \begin{block}{PRM (Process Reward Models)}
              Работают на \textbf{уровне шагов}, присваивая награду каждому промежуточному шагу, обусловленному задачей и всеми предыдущими шагами. \\
              \vspace{0.5em}
              \footnotesize \textit{Стандартный подход для верификации рассуждений благодаря более точной гранулярности.}
          \end{block}
      \end{column}
  \end{columns}
  
  \vspace{0.5em}
  \begin{block}{}
      \textbf{Step-level супервизия} позволяет PRM идентифицировать галлюцинированные или некорректные шаги в цепочке рассуждений. Благодаря этой более точной гранулярности, PRM стали стандартом для верификации.
  \end{block}
\end{frame}

% Слайд 4
\begin{frame}
  \frametitle{Ограничения: Проблемы текущих датасетов}
  \footnotesize
  
  \begin{columns}[T]
      \begin{column}{0.48\textwidth}
          \begin{block}{1. Гранулярность супервизии}
              Супервизия почти исключительно предоставляется на \textbf{уровне шагов}. \\
              На практике галлюцинации часто происходят на \textbf{более мелкой гранулярности}:
              \begin{itemize}
                  \item Арифметическая ошибка внутри шага
                  \item Необоснованное преобразование внутри иначе корректного шага
              \end{itemize}
          \end{block}
      \end{column}
      
      \begin{column}{0.48\textwidth}
          \begin{block}{2. Масштаб датасетов}
              Многие датасеты типично большие — сотни тысяч примеров. \\
              Это создаёт проблемы:
              \begin{itemize}
                  \item Построение датасетов становится дорогостоящим
                  \item Обучение PRM на них требует значительных вычислительных ресурсов
              \end{itemize}
          \end{block}
      \end{column}
  \end{columns}

  \vspace{0.5em}
  \begin{alertblock}{Необходимость}
      Методология для генерации span-level галлюцинационных аннотаций в математических рассуждениях с автоматизацией процесса.
  \end{alertblock}
\end{frame}

% Слайд 5
\begin{frame}
  \frametitle{Решение: SHARP-Math}
  \framesubtitle{Высококачественный датасет для детекции галлюцинаций}
  \footnotesize

  \begin{columns}[T]
      \begin{column}{0.5\textwidth}
          \begin{block}{Ключевое отличие}
              В отличие от существующих ресурсов (с метками на уровне шагов), SHARP-Math предоставляет \textbf{детальные аннотации на уровне span (фрагментов)}, сгенерированные через автоматизированный pipeline.
          \end{block}
          \begin{block}{Методология}
              Синтезирует \textbf{реалистичные траектории рассуждений} с точной разметкой ошибок.
          \end{block}
          \begin{itemize}
              \item \textbf{11K} примеров
              \item \textbf{88\%} с ошибками
              \item \textbf{12\%} полностью корректных
          \end{itemize}
      \end{column}
      
      \begin{column}{0.48\textwidth}
          \begin{block}{Достижения}
              \begin{itemize}
                  \item \textbf{Первый датасет} специально для span-level детекции галлюцинаций (автоматическая генерация).
                  \item \textbf{State-of-the-art} производительность на ProcessBench (PRM 1.5B достигает 56.3 F1).
                  \item \textbf{В 1-2 порядка меньше} данных по сравнению с открытыми датасетами.
              \end{itemize}
          \end{block}
      \end{column}
  \end{columns}
\end{frame}

% Слайд 6
\begin{frame}
  \frametitle{Методология: Трёхэтапный pipeline}
  \footnotesize
  
  \begin{columns}[T]
      \begin{column}{0.31\textwidth}
          \begin{block}{1. Выбор задач}
              Задачи и эталонные решения из датасета \textbf{MATH}. Выбран из-за сложных задач и широкого охвата математических дисциплин.
          \end{block}
      \end{column}
      
      \begin{column}{0.34\textwidth}
          \begin{block}{2. Генерация решений}
              Использование \textbf{Mistral-7B-Instruct} для генерации траекторий рассуждений. Модель создаёт связные пошаговые решения с естественным распределением ошибок.
          \end{block}
      \end{column}

      \begin{column}{0.31\textwidth}
          \begin{block}{3. Span-аннотации}
              Детекция с помощью \textbf{gpt-oss-120B} как автоматического аннотатора. Сравнивает решение с эталоном и оборачивает фрагменты в теги.
          \end{block}
      \end{column}
  \end{columns}
  
  \vspace{0.5em}
  \begin{block}{Фильтрация и очистка}
      \small Применяются эвристические фильтры: исключение примеров со спанами длиннее 110 символов, более 35 спанов, или с malformed генерациями. 
      После фильтрации: \textbf{11,678 высококачественных примеров}.
  \end{block}
\end{frame}

% Слайд 7
\begin{frame}
  \frametitle{Сравнение с существующими датасетами}
  \scriptsize
  
  \begin{table}[h]
  \scriptsize
  \begin{tabularx}{\textwidth}{Y Y Y Y Y}
      \toprule
      \textbf{Датасет} & \textbf{Разметка} & \textbf{Размер} & \textbf{Доступн.} & \textbf{Гранул.} \\
      \midrule
      PRM800K & Ручная & 415K & Открытый & Step-level \\
      QwenPRM & Авто & 1,500K & Закрытый & Step-level \\
      Math-Shepherd & Авто & 445K & Открытый & Step-level \\
      RLHFlow-PRM & Авто & 253K & Открытый & Step-level \\
      \midrule
      \textbf{SHARP-Math (ours)} & \textbf{Авто} & \textbf{11K} & \textbf{Открытый} & \textbf{Span-level} \\
      \bottomrule
  \end{tabularx}
  \end{table}

  \vspace{0.5em}
  \begin{columns}[T]
      \begin{column}{0.3\textwidth}
          \scriptsize\textbf{Компактность:} \\ В 10-50 раз меньше данных при более плотной супервизии
      \end{column}
      \begin{column}{0.33\textwidth}
          \scriptsize\textbf{Точность:} \\ Span-level аннотации точнее локализуют ошибки
      \end{column}
      \begin{column}{0.3\textwidth}
          \scriptsize\textbf{Доступность:} \\ Полностью открытый датасет для сообщества
      \end{column}
  \end{columns}
\end{frame}

% Слайд 8
\begin{frame}
  \frametitle{Результаты: Производительность на ProcessBench}
  \small
  
  \begin{columns}[T]
      \begin{column}{0.48\textwidth}
          \begin{block}{7B Модель (Qwen2.5-Math-7B)}
              \textbf{SHARP-Math-Span:} 67.0 F1 \\
              PRM800K baseline: 56.5 F1 \\
              \textit{+10.4 пунктов улучшения}
          \end{block}
          
          \begin{block}{1.5B Модель (Qwen2.5-Math-1.5B)}
              \textbf{SHARP-Math-Span:} 60.0 F1 \\
              7B PRM800K baseline: 56.5 F1 \\
              \textit{+3.5 пункта улучшения}
          \end{block}
      \end{column}
      
      \begin{column}{0.48\textwidth}
          \begin{block}{Сравнение с другими моделями (7B)}
              \scriptsize
              \begin{tabularx}{\linewidth}{Y r}
                  Qwen2.5-Math-PRM-7B & 73.5 F1 \\
                  R-PRM-7B-DPO & 70.4 F1 \\
                  PathFinder-PRM-7B & 69.5 F1 \\
                  \textbf{SHARP-Math-Span (ours)} & \textbf{67.0 F1} \\
                  Qwen2.5-Math-7B-PRM800K & 56.5 F1 \\
                  Skywork-PRM-7B & 42.1 F1 \\
              \end{tabularx}
          \end{block}
          \footnotesize \textit{Конкурентная обобщаемость на всех датасетах: GSM8K, MATH, OlympiadBench, Omni-MATH.}
      \end{column}
  \end{columns}
\end{frame}

% Слайд 9
\begin{frame}
  \frametitle{Эффективность данных и масштабирование}
  
  \begin{columns}[T]
      \begin{column}{0.48\textwidth}
          \begin{block}{Размер датасета}
              \textbf{SHARP-Math:} 11K (Наш подход) \\
              R-PRM-7B: 285K-553K \\
              RetrievalPRM-7B: 445K \\
              \vspace{0.5em}
              \textit{В 10-50 раз меньше данных}
          \end{block}
          \begin{block}{Performance-per-sample}
              Плотная супервизия в SHARP-Math обеспечивает \textbf{более высокое качество на единицу данных} по сравнению с традиционными датасетами уровня шагов.
          \end{block}
      \end{column}
      
      \begin{column}{0.48\textwidth}
          \begin{block}{Кривая обучения (F1 vs Токены)}
              % Заглушка для графика. В оригинале здесь график со сравнением скорости достижения нужного F1 Score.
              \centering
              %\includegraphics[width=\textwidth]{learning_curve_placeholder.pdf}
              \textit{[График: SHARP-Math быстрее достигает высокого F1, тогда как PRM800K имеет более медленное улучшение]}
              
              \vspace{0.5em}
              \footnotesize
              \begin{itemize}
                  \item \textbf{SHARP-Math:} F1 $\sim$44.6 после 0.45M токенов
                  \item \textbf{PRM800K:} F1 > 60 при $\sim$2.0M токенов
              \end{itemize}
          \end{block}
      \end{column}
  \end{columns}
\end{frame}

% Слайд 10
\begin{frame}
  \frametitle{Абляционное исследование гранулярности}
  
  \small Для проверки вклада источника данных и гранулярности аннотаций проводилось обучение одних и тех же PRM backbone при идентичных гиперпараметрах на двух контролируемых датасетах: SHARP-Math-Span и SHARP-Math-Step.
  
  \vspace{0.5em}
  \begin{columns}[T]
      \begin{column}{0.48\textwidth}
          \begin{block}{7B Backbone}
              \textbf{SHARP-Math-Span:} 67.0 F1 \textit{(Лучший результат)} \\
              SHARP-Math-Step: 65.8 F1 \textit{(Базовый уровень)} \\
              \vspace{0.5em}
              \textbf{Вывод: Span лучше для 7B}
          \end{block}
      \end{column}
      
      \begin{column}{0.48\textwidth}
          \begin{block}{1.5B Backbone}
              \textbf{SHARP-Math-Span:} 60.0 F1 \textit{(Лучший результат)} \\
              SHARP-Math-Step: 59.2 F1 \textit{(Базовый уровень)} \\
              \vspace{0.5em}
              \textbf{Вывод: Span лучше для 1.5B}
          \end{block}
      \end{column}
  \end{columns}

  \vspace{0.5em}
  \begin{block}{Ключевой вывод}
      \footnotesize Эффект span-аннотаций не зависит от \textbf{ёмкости модели}. Span-аннотации обеспечивают более богатый сигнал обучения, а модели обученные на таких аннотациях превосходят модели обученные на step-аннотациях в обоих размерах.
  \end{block}
\end{frame}

% Слайд 11
\begin{frame}
  \frametitle{Ограничения и этические соображения}
  \footnotesize
  
  \begin{columns}[T]
      \begin{column}{0.48\textwidth}
          \begin{block}{Ограничения датасета}
              \begin{enumerate}
                  \item Охватывает \textbf{только задачи из MATH} — ограниченная область.
                  \item Сгенерированы \textbf{одной небольшой моделью} (вносит смещение).
                  \item Требует \textbf{дополнительной валидации} с измерением меж-аннотаторного согласия.
              \end{enumerate}
          \end{block}
          
          \begin{block}{Планы по улучшению}
              Расширение датасета для дополнительных задач и интеграция \textbf{нескольких моделей} для генерации галлюцинаций.
          \end{block}
      \end{column}
      
      \begin{column}{0.48\textwidth}
          \begin{block}{Этические соображения}
              \textbf{Риск чрезмерного доверия:} Пользователи могут чрезмерно полагаться на детекцию ошибок. \\
              \textbf{Митигация:} Чёткая документация \textbf{ограничений и предполагаемого исследовательского использования.}
          \end{block}
          
          \begin{block}{Приватность данных}
              \begin{itemize}
                  \item Не содержит персональной информации
                  \item Данные полностью сгенерированы моделью
                  \item Нет данных от человеческих пользователей
              \end{itemize}
          \end{block}
      \end{column}
  \end{columns}
\end{frame}

% Слайд 12
\begin{frame}
  \frametitle{Ключевые выводы: Итоги исследования}
  \footnotesize
  
  \begin{columns}[T]
      \begin{column}{0.48\textwidth}
          \begin{block}{Первый span-level датасет}
              \textbf{SHARP-Math} — первый открытый датасет с детальными span-level аннотациями для математических ошибок рассуждений.
          \end{block}
          \begin{block}{Превосходная производительность}
              Обучение PRM на плотном span-level сигнале даёт \textbf{превосходное качество} даже при небольшом количестве данных.
          \end{block}
      \end{column}
      
      \begin{column}{0.48\textwidth}
          \begin{block}{Автоматизация процесса}
              Автоматизация аннотации решает проблемы \textbf{масштабируемости и гранулярности} существующих датасетов.
          \end{block}
          \begin{block}{Открытость для сообщества}
              Полностью открытый датасет для \textbf{исследовательского сообщества} и дальнейшего развития.
          \end{block}
      \end{column}
  \end{columns}
  
  \vspace{0.5em}
  \begin{center}
      \scriptsize\textit{\textbf{SHARP-Math} создаёт основу для data-efficient верификаторов рассуждений и более надёжных математических LLM.}
  \end{center}
\end{frame}

% Финальный слайд
\begin{frame}
\begin{center}
{\Large Спасибо за~внимание!}
\end{center}
\end{frame}

\end{document}